\documentclass[a4paper,11pt]{article}

% --- Packages ---
\usepackage[utf8]{inputenc}
\usepackage[T1]{fontenc}
\usepackage[french]{babel}
\usepackage{geometry}
\usepackage{graphicx}
\usepackage{booktabs}
\usepackage{longtable}
\usepackage{xcolor}
\usepackage{hyperref}
\usepackage{fancyhdr}
\usepackage{titlesec}
\usepackage{enumitem}
\usepackage{amsmath}

\geometry{margin=2.5cm}

% --- Couleurs ---
\definecolor{bleuaudit}{HTML}{3498DB}
\definecolor{vertok}{HTML}{27AE60}
\definecolor{orangepreoc}{HTML}{F39C12}
\definecolor{rougecrit}{HTML}{E74C3C}
\definecolor{grisfonce}{HTML}{2C3E50}

% --- En-tete/pied de page ---
\pagestyle{fancy}
\fancyhf{}
\fancyhead[L]{\small\textcolor{grisfonce}{Rapport d'audit IA --- \VAR{audit_name}}}
\fancyhead[R]{\small\textcolor{grisfonce}{\VAR{date}}}
\fancyfoot[C]{\thepage}
\fancyfoot[R]{\small\textcolor{grisfonce}{Confidentiel}}

% --- Titre ---
\titleformat{\section}{\Large\bfseries\color{grisfonce}}{}{0em}{}[\vspace{-0.5em}\rule{\textwidth}{0.5pt}]
\titleformat{\subsection}{\large\bfseries\color{bleuaudit}}{}{0em}{}

% --- Jinja2 : les variables sont entourees de \VAR{} et les blocs de \BLOCK{} ---
% Compilation : d'abord rendre via Jinja2, puis compiler le .tex resultant

\begin{document}

% === PAGE DE GARDE ===
\begin{titlepage}
\centering
\vspace*{3cm}

{\Huge\bfseries\textcolor{grisfonce}{Rapport d'Audit IA}}\\[0.5cm]
{\LARGE\textcolor{bleuaudit}{\VAR{audit_name}}}\\[2cm]

\begin{tabular}{ll}
\textbf{Client}   & \VAR{client_name} \\
\textbf{Systeme}  & \VAR{system_name} \\
\textbf{Auditeur} & \VAR{auditor_name} \\
\textbf{Date}     & \VAR{date} \\
\textbf{Version}  & \VAR{version} \\
\textbf{Statut}   & \VAR{status} \\
\end{tabular}

\vfill

{\large\textcolor{grisfonce}{Methodologie d'audit LLM v1.0}}\\
{\small Hanen Mizouni --- IA au feminin}

\vspace{1cm}
{\footnotesize\textcolor{rougecrit}{CONFIDENTIEL --- Ne pas diffuser sans autorisation}}
\end{titlepage}

% === TABLE DES MATIERES ===
\tableofcontents
\newpage

% === 1. SYNTHESE EXECUTIVE ===
\section{Synthese executive}

Ce rapport presente les resultats de l'audit du systeme IA \textbf{\VAR{system_name}},
realise du \VAR{date_start} au \VAR{date_end} selon la methodologie d'audit LLM v1.0.

\subsection{Score global}

\begin{center}
\begin{tabular}{lccc}
\toprule
\textbf{Dimension} & \textbf{Score} & \textbf{Grade} & \textbf{Poids} \\
\midrule
Donnees (etape 2) & \VAR{score_data}/100 & \VAR{grade_data} & 20\% \\
Fairness (etape 3) & \VAR{score_fairness}/100 & \VAR{grade_fairness} & 30\% \\
Red-teaming (etape 4) & \VAR{score_redteam}/100 & \VAR{grade_redteam} & 25\% \\
Robustesse (etape 5) & \VAR{score_robustness}/100 & \VAR{grade_robustness} & 15\% \\
Conformite & \VAR{score_compliance}/100 & \VAR{grade_compliance} & 10\% \\
\midrule
\textbf{GLOBAL} & \textbf{\VAR{score_global}/100} & \textbf{\VAR{grade_global}} & 100\% \\
\bottomrule
\end{tabular}
\end{center}

\subsection{Top 3 des risques identifies}

\begin{enumerate}
\BLOCK{for risk in top_risks}
    \item \textbf{\VAR{risk.title}} --- Severite : \VAR{risk.severity}\\
    \VAR{risk.description}
\BLOCK{endfor}
\end{enumerate}

\subsection{Top 3 des recommandations prioritaires}

\begin{enumerate}
\BLOCK{for rec in top_recommendations}
    \item \textbf{\VAR{rec.title}} --- Effort : \VAR{rec.effort}\\
    \VAR{rec.description}
\BLOCK{endfor}
\end{enumerate}

\newpage

% === 2. CONTEXTE ET METHODOLOGIE ===
\section{Contexte et methodologie}

\subsection{Description du systeme}
\VAR{system_description}

\subsection{Perimetre de l'audit}
\VAR{audit_scope}

\subsection{Methodologie appliquee}
Cet audit suit la \textbf{Methodologie d'audit LLM v1.0} (Mizouni, 2026),
structuree en 7 etapes sur 5 jours. Les cadres de reference sont :

\begin{itemize}
    \item EU AI Act (Reglement 2024/1689) --- Annexes IV et XI
    \item NIST AI Risk Management Framework (AI 100-1)
    \item ISO/IEC 42001 --- Systeme de management de l'IA
    \item Recommandations CNIL sur l'IA
\end{itemize}

\subsection{Limites de l'audit}
\VAR{audit_limitations}

% === 3. ANALYSE DES DONNEES ===
\section{Analyse des donnees (etape 2)}

\VAR{data_analysis_content}

% === 4. RESULTATS DE FAIRNESS ===
\section{Resultats de fairness (etape 3)}

\VAR{fairness_content}

% === 5. RESULTATS DU RED-TEAMING ===
\section{Red-teaming et tests adversariaux (etape 4)}

\VAR{redteam_content}

% === 6. RESULTATS DE ROBUSTESSE ===
\section{Robustesse et stabilite (etape 5)}

\VAR{robustness_content}

% === 7. PLAN DE REMEDIATION ===
\section{Plan de remediation (etape 6)}

\VAR{remediation_content}

% === 8. CONFORMITE AI ACT ===
\section{Conformite AI Act}

\subsection{Classification du systeme}
\begin{itemize}
    \item \textbf{Niveau de risque} : \VAR{risk_level}
    \item \textbf{Justification} : \VAR{risk_justification}
\end{itemize}

\subsection{Obligations applicables}
\begin{longtable}{p{6cm}p{3cm}p{4cm}}
\toprule
\textbf{Obligation} & \textbf{Statut} & \textbf{Action requise} \\
\midrule
\BLOCK{for obligation in obligations}
\VAR{obligation.name} & \VAR{obligation.status} & \VAR{obligation.action} \\
\BLOCK{endfor}
\bottomrule
\end{longtable}

% === 9. ANNEXES ===
\section{Annexes}

\subsection{A. Methodologie detaillee}
Voir documentation complete : \url{https://github.com/[repo]/audit-llm-methodology}

\subsection{B. Catalogue de prompts utilises}
\VAR{prompts_catalog_summary}

\subsection{C. Tableaux de metriques detailles}
\VAR{detailed_metrics_tables}

\subsection{D. Bibliographie}
\begin{itemize}
    \item Barocas, S. \& Selbst, A. (2016). Big Data's Disparate Impact. \textit{California Law Review}.
    \item Chouldechova, A. (2017). Fair prediction with disparate impact. \textit{Big Data}.
    \item Kleinberg, J., Mullainathan, S. \& Raghavan, M. (2016). Inherent Trade-Offs in the Fair Determination of Risk Scores. \textit{ITCS}.
    \item Mitchell, M. et al. (2019). Model Cards for Model Reporting. \textit{FAT*}.
    \item NIST (2023). AI Risk Management Framework. AI 100-1.
    \item Reglement (UE) 2024/1689 --- AI Act.
\end{itemize}

\vfill

\begin{center}
\small\textcolor{grisfonce}{
Rapport genere selon la Methodologie d'audit LLM v1.0\\
Hanen Mizouni --- IA au feminin\\
\VAR{date}
}
\end{center}

\end{document}
